\lecture{3}{Graph Neural Networks I}{Tuesday, September 29, 2026}\label{ch:lec03}

% ---------------------------------------------------------------------
%  Local TikZ styles (prefix lec03) and a macro for the 6-node example
%  graph of the slides (edges A-B, A-C, A-D, B-C, C-E, C-F, E-F).
%  Node colour is tied to the POSITION, so it follows a node when it is
%  relabelled; #3..#8 are the labels of the six positions.
% ---------------------------------------------------------------------
\tikzset{
  lec03cA/.style={fill=Tan!80!Brown},
  lec03cB/.style={fill=BrickRed!65},
  lec03cC/.style={fill=ForestGreen!55},
  lec03cD/.style={fill=RoyalBlue!50},
  lec03cE/.style={fill=Orchid!65},
  lec03cF/.style={fill=WildStrawberry!50},
  lec03s/.style={circle, draw=black!60, minimum size=3.6mm, inner sep=0pt,
                 font=\tiny\sffamily},
  lec03m/.style={circle, draw=black!60, thick, minimum size=4.8mm,
                 inner sep=0pt, font=\scriptsize\sffamily},
  lec03bxl/.style={draw=black!60, fill=black!12, minimum size=2.6mm,
                   inner sep=0pt},
  lec03bxd/.style={draw=black!70, fill=black!60, minimum size=1.8mm,
                   inner sep=0pt},
  lec03msg/.style={black!45, -{Stealth[length=1.1mm]}},
}
\providecommand{\lecIIIsixgraph}[8]{%
  \node[#1, lec03cA] (#2a) at (0.9,1.25) {#3};
  \node[#1, lec03cB] (#2b) at (2.5,1.9)  {#4};
  \node[#1, lec03cC] (#2c) at (2.9,1.0)  {#5};
  \node[#1, lec03cD] (#2d) at (0.0,0.0)  {#6};
  \node[#1, lec03cE] (#2e) at (1.9,0.1)  {#7};
  \node[#1, lec03cF] (#2f) at (3.6,0.35) {#8};
  \draw[edge] (#2a)--(#2b) (#2a)--(#2c) (#2a)--(#2d) (#2b)--(#2c)
              (#2c)--(#2e) (#2c)--(#2f) (#2e)--(#2f);
}

\textbf{Outline.} Recap $\to$ deep learning for graphs (permutation
invariance / equivariance) $\to$ graph convolutional networks (GCN) $\to$ GNNs
subsume CNNs (and transformers) $\to$ background: basics of deep learning (the
deck's appendix, slides 73--88, kept at the end in deck order).

% =====================================================================
\section{From shallow to deep graph encoders}\label{sec:lec03-intro}

\subsection{Recap: node embeddings and shallow encoders}\label{sec:lec03-recap}
\slides{3--6}

\begin{itemize}
  \item \textbf{Intuition:} map nodes to $d$-dimensional embeddings such that
    similar nodes in the graph are embedded close together:
    $f(\text{input graph}) = $ 2-D node embeddings (figure: disease network;
    related diseases cluster, e.g.\ around thyroid disease, viral meningitis,
    pancreatitis, uterine polyp). \textbf{How to \emph{learn} the mapping $f$?}
  \item \textbf{Goal:} encode nodes ($\ENC(u)\to\vz_u$, $\ENC(v)\to\vz_v$ in the
    $d$-dim embedding space) so that
\end{itemize}
\begin{equation}\label{eq:lec03-recap}
  \similarity(u,v) \approx \vz_v\tp\vz_u
  \qquad\text{(both $\similarity$ and $\ENC$: ``need to define!'')}
\end{equation}
\begin{itemize}
  \item \term{Shallow encoding}: the encoder is just an \textbf{embedding
    lookup}; embedding matrix $\mZ$ with one column per node (column =
    embedding vector of that node), number of rows = dimension/size of the
    embeddings.
  \item \textbf{Limitations} of shallow embedding methods:
    \begin{itemize}
      \item \textbf{$\bigO(|V|d)$ parameters are needed}: no sharing of
        parameters between nodes; every node has its own unique embedding.
      \item \textbf{Inherently ``transductive''}: cannot generate embeddings for
        nodes not seen during training.
      \item \textbf{Do not incorporate node features}: nodes in many graphs have
        features that we can and should leverage.
    \end{itemize}
\end{itemize}

\subsection{Deep graph encoders; why graphs are hard}\label{sec:lec03-deep}
\slides{7--13}

\begin{itemize}
  \item \textbf{Today:} deep learning methods based on \term{graph neural
    networks} (GNNs):
\end{itemize}
\begin{equation}\label{eq:lec03-enc}
  \ENC(v) = \text{multiple layers of non-linear transformations based on graph structure}
\end{equation}
\begin{itemize}
  \item \textbf{Note:} all these deep encoders can be \textbf{combined with the
    node similarity functions} defined in Lecture~2.
  \item Architecture (slide 8): graph convolutions $\to$ activation function
    $\to$ regularisation (e.g.\ dropout) $\to$ graph convolutions $\to\cdots$
    (each layer: one vector per node). \textbf{Output:} node embeddings; we can
    also embed subgraphs and graphs.
  \item \textbf{Tasks} we will be able to solve: \textbf{node classification}
    (predict the type of a given node); \textbf{link prediction} (predict
    whether two nodes are linked); \textbf{community detection} (identify
    densely linked clusters of nodes); \textbf{network similarity} (how
    similar two (sub)networks are).
  \item \textbf{Modern ML toolbox:} images = grids (CNN layers: patterns of
    local contrast $\to$ face features $\to$ faces); text/speech = sequences
    (chain of nodes; RNN cells). \textbf{Modern deep learning toolbox is designed for
    simple sequences \& grids.}
  \item \textbf{Why is it hard? Networks are far more complex:} arbitrary size
    and complex topological structure (i.e.\ no spatial locality like grids);
    no fixed node ordering or reference point; often dynamic and have
    multimodal features.
  \item Outline (slide 12): deep learning for graphs; graph convolutional
    networks; GNNs subsume CNNs.
  \item \textbf{Background: basics of deep learning} (slide 13; details in
    \cref{sec:lec03-dl}): loss function $\Loss$,
\end{itemize}
\begin{equation}\label{eq:lec03-dl-recipe}
  \min_{\Theta}\ \Loss\bigl(\vy, f_\Theta(\vx)\bigr)
\end{equation}
\begin{itemize}
  \item $f$ can be a simple linear layer, an MLP, or other neural networks
    (e.g.\ a GNN later). Sample a minibatch of input data $\vx$.
    \textbf{Forward propagation:} compute $\Loss$ given $\vx$.
    \textbf{Back-propagation:} obtain the gradient $\nabla_\Theta\Loss$ using a
    chain rule. Use \textbf{stochastic gradient descent (SGD)} to optimise
    $\Loss$ for weights $\Theta$ over many iterations.
\end{itemize}

% =====================================================================
\section{Deep learning for graphs}\label{sec:lec03-dlg}
\slides{14--34}

\textbf{Content} (slide 15): \textbf{local network neighbourhoods} (describe
aggregation strategies; define computation graphs) and \textbf{stacking
multiple layers} (describe the model, parameters, training; how to fit the
model; simple example for unsupervised and supervised training).

\subsection{Setup and a naive approach}\label{sec:lec03-setup}
\slides{16--17}

\begin{itemize}
  \item Graph $G$: $V$ = \textbf{vertex set}; $\mA$ = \textbf{adjacency
    matrix} (assume binary); $\mX\in\R^{|V|\times m}$ = matrix of
    \textbf{node features}; $v$ = a node in $V$; $\Nb(v)$ = the set of
    neighbours of $v$.
  \item \textbf{Node features:} social networks: user profile, user image;
    biological networks: gene expression profiles, gene functional information.
    \textbf{No node features} in the dataset: indicator vectors (one-hot
    encoding of a node), or a vector of constant 1: $[1,1,\dots,1]$.
  \item \textbf{Naive approach:} join adjacency matrix and features $[\mA,\mX]$
    and feed them (one row per node) into a deep neural net
    (\cref{fig:lec03-naive}).
  \item \textbf{Issues:} $\bigO(|V|)$ parameters; not applicable to graphs of
    different sizes; sensitive to node ordering.
\end{itemize}

\begin{figure}[htbp]
  \centering
  \begin{tikzpicture}[scale=0.85, transform shape]
    \begin{scope}[every node/.style={vertex, minimum size=5mm, font=\scriptsize}]
      \node (A) at (0,1.1) {A}; \node (B) at (1.2,1.1) {B};
      \node (C) at (0,0) {C};   \node (D) at (1.2,0) {D};
      \node (E) at (2.0,0.55) {E};
    \end{scope}
    \draw[edge] (A)--(B) (A)--(C) (A)--(D) (B)--(D) (B)--(E) (C)--(D) (D)--(E);
    \matrix (m) at (4.7,0.55) [matrix of math nodes,
        nodes={minimum width=4.6mm, minimum height=3.9mm, inner sep=0pt, font=\scriptsize}] {
        & A & B & C & D & E & f_1 & f_2 \\
      A & 0 & 1 & 1 & 1 & 0 & 1 & 0 \\
      B & 1 & 0 & 0 & 1 & 1 & 0 & 0 \\
      C & 1 & 0 & 0 & 1 & 0 & 0 & 1 \\
      D & 1 & 1 & 1 & 0 & 1 & 1 & 1 \\
      E & 0 & 1 & 0 & 1 & 0 & 1 & 0 \\
    };
    \draw[black!50] (m-1-6.north east) -- (m-6-6.south east);
    \draw[SkyBlue!90!black, thick, rounded corners=1pt] (m-3-2.north west) rectangle (m-3-8.south east);
    \node[font=\scriptsize, above] at (m-1-4.north) {$\mA$};
    \node[font=\scriptsize, above] at ($(m-1-7.north)!0.5!(m-1-8.north)$) {$\mX$ (feat.)};
    \begin{scope}[shift={(8.0,0)}, every node/.style={circle, draw, minimum size=2.6mm, inner sep=0pt}]
      \foreach \i in {0,...,6} \node (p\i) at (0,1.35-0.27*\i) {};
      \foreach \i in {0,...,5} \node (q\i) at (0.9,1.3-0.27*\i) {};
      \foreach \i in {0,...,5} \node (r\i) at (1.8,1.3-0.27*\i) {};
      \foreach \i in {0,...,2} \node (s\i) at (2.7,0.85-0.27*\i) {};
      \foreach \i in {0,...,6} \foreach \j in {0,...,5} \draw[black!25] (p\i) -- (q\j);
      \foreach \i in {0,...,5} \foreach \j in {0,...,5} \draw[black!25] (q\i) -- (r\j);
      \foreach \i in {0,...,5} \foreach \j in {0,...,2} \draw[black!25] (r\i) -- (s\j);
    \end{scope}
    \draw[SkyBlue!90!black, thick, rounded corners=1pt] (7.82,-0.43) rectangle (8.18,1.53);
    \draw[arrow] ([xshift=4pt]m-3-8.east) -- (7.72,0.55 |- m-3-8.east);
    \node[font=\Large\bfseries] at (11.2,0.55) {?};
  \end{tikzpicture}
  \caption{Naive approach (slides 17, 33): row $v$ of $[\mA,\mX]$ (here $B$)
  is the input layer of an MLP.}
  \label{fig:lec03-naive}
\end{figure}

\subsection{Convolutions? Permutation invariance}\label{sec:lec03-perm-inv}
\slides{18--25}

\begin{itemize}
  \item \textbf{Idea: convolutional networks.} CNN on an image: a $3\times3$
    window slides over the pixel grid; input $\to$ convolutions $\to$ feature
    maps $\to$ subsampling $\to$ convolutions $\to$ subsampling $\to$ fully
    connected $\to$ output. \textbf{Goal:} generalise convolutions beyond
    simple lattices; leverage node features/attributes (e.g.\ text, images).
  \item \textbf{Real-world graphs:} there is \textbf{no fixed notion of
    locality or sliding window} on the graph; graphs are \textbf{permutation
    invariant}.
  \item Imagine we want to \textbf{embed an entire graph}.
    \textbf{Observation: a graph does not have a canonical ordering of its
    nodes}; we can have many different node orderings of the same graph.
    \textbf{What we want:} if we learn an embedding function over a graph,
    we should get the same result (the same embedding) regardless of how the
    nodes are numbered.
  \item Running example (\cref{fig:lec03-orders}): order~1 $(\mA_1,\mX_1)$ and
    order~2 $(\mA_2,\mX_2)$ of the same graph (order~2 renames $A,\dots,F$ of
    order~1 to $E,F,D,C,B,A$). \textbf{Learned graph representation should be
    the same for order~1 and order~2.}
  \item Formally: learn $f$ that maps a graph $G=(\mA,\mX)$ to a vector in
    $\R^d$ (a $d$-dim embedding); then we want
    $f(\mA_1,\mX_1) = f(\mA_2,\mX_2)$. If $f(\mA_i,\mX_i)=f(\mA_j,\mX_j)$ for
    any orderings $i$ and $j$, $f$ is a \textbf{permutation invariant
    function} (a graph with $|V|$ nodes has $|V|!$ different orderings).
\end{itemize}

\begin{definition}[Permutation invariance]\label{def:lec03-perm-inv}
  A graph function $f:\R^{|V|\times|V|}\times\R^{|V|\times m}\to\R^{d}$ is
  \term{permutation invariant} if
  \begin{equation}\label{eq:lec03-perm-inv}
    f(\mA,\mX) = f(\mP\mA\mP\tp,\mP\mX)\qquad\text{for any permutation } \mP .
  \end{equation}
  $m$: each node has an $m$-dim feature vector; $d$: output embedding
  dimensionality of the graph $G=(\mA,\mX)$; permutation $\mP$: a shuffle of
  the node order, e.g.\ $(A,B,C)\to(B,C,A)$.
\end{definition}

\begin{figure}[htbp]
  \centering
  \begin{tikzpicture}[scale=0.8, transform shape, lbl/.style={font=\scriptsize\sffamily, inner sep=1pt}]
    % ================= ORDER 1 =================
    \node[font=\small\bfseries, cautionred, anchor=west] at (-0.5,2.55) {Order 1: $(\mA_1,\mX_1)$};
    \lecIIIsixgraph{vertex}{o}{A}{B}{C}{D}{E}{F}
    \node[font=\scriptsize] at (5.15,2.25) {features $\mX_1$};
    \foreach \l/\c [count=\r] in {A/A,B/B,C/C,D/D,E/E,F/F} {
      \node[lbl] at (4.35,1.95-0.32*\r+0.16) {\l};
      \fill[lec03c\c, rounded corners=1.2mm] (4.55,1.95-0.32*\r+0.04) rectangle (5.75,1.95-0.32*\r+0.28);
    }
    \begin{scope}[shift={(6.5,0.03)}]
      \foreach \r/\c in {1/2,1/3,1/4,2/1,2/3,3/1,3/2,3/5,3/6,4/1,5/3,5/6,6/3,6/5}
        \fill[SkyBlue!80] ({0.32*(\c-1)},{1.92-0.32*\r}) rectangle ++(0.32,0.32);
      \draw[black!45] (0,0) grid[step=0.32] (1.92,1.92);
      \foreach \l [count=\i] in {A,B,C,D,E,F} {
        \node[lbl] at ({0.32*\i-0.16},2.05) {\l};
        \node[lbl] at (-0.14,{1.92-0.32*\i+0.16}) {\l};
      }
      \node[font=\scriptsize] at (0.96,2.4) {adjacency $\mA_1$};
    \end{scope}
    \begin{scope}[shift={(9.3,0.03)}]
      \foreach \l/\c [count=\r] in {A/A,B/B,C/C,D/D,E/E,F/F} {
        \node[lbl] at (-0.16,{1.92-0.32*\r+0.16}) {\l};
        \fill[lec03c\c] (0,{1.92-0.32*\r}) rectangle ++(0.64,0.32);
      }
      \draw[black!45] (0,0) grid[step=0.32] (0.64,1.92);
      \draw[very thick, Goldenrod!70!black] (-0.05,1.55) rectangle (0.69,1.97);
      \draw[thick, dashed] (-0.05,0.59) rectangle (0.69,1.01);
      \node[font=\scriptsize] at (0.32,2.4) {$f(\mA_1,\mX_1)$};
    \end{scope}
    % ================= ORDER 2 =================
    \begin{scope}[shift={(0,-3.3)}]
      \node[font=\small\bfseries, asideteal, anchor=west] at (-0.5,2.55) {Order 2: $(\mA_2,\mX_2)$};
      \lecIIIsixgraph{vertex}{t}{E}{F}{D}{C}{B}{A}
      \node[font=\scriptsize] at (5.15,2.25) {features $\mX_2$};
      \foreach \l/\c [count=\r] in {A/F,B/E,C/D,D/C,E/A,F/B} {
        \node[lbl] at (4.35,1.95-0.32*\r+0.16) {\l};
        \fill[lec03c\c, rounded corners=1.2mm] (4.55,1.95-0.32*\r+0.04) rectangle (5.75,1.95-0.32*\r+0.28);
      }
      \begin{scope}[shift={(6.5,0.03)}]
        \foreach \r/\c in {1/2,1/4,2/1,2/4,3/5,4/1,4/2,4/5,4/6,5/3,5/4,5/6,6/4,6/5}
          \fill[SkyBlue!80] ({0.32*(\c-1)},{1.92-0.32*\r}) rectangle ++(0.32,0.32);
        \draw[black!45] (0,0) grid[step=0.32] (1.92,1.92);
        \foreach \l [count=\i] in {A,B,C,D,E,F} {
          \node[lbl] at ({0.32*\i-0.16},2.05) {\l};
          \node[lbl] at (-0.14,{1.92-0.32*\i+0.16}) {\l};
        }
        \node[font=\scriptsize] at (0.96,2.4) {adjacency $\mA_2$};
      \end{scope}
      \begin{scope}[shift={(9.3,0.03)}]
        \foreach \l/\c [count=\r] in {A/F,B/E,C/D,D/C,E/A,F/B} {
          \node[lbl] at (-0.16,{1.92-0.32*\r+0.16}) {\l};
          \fill[lec03c\c] (0,{1.92-0.32*\r}) rectangle ++(0.64,0.32);
        }
        \draw[black!45] (0,0) grid[step=0.32] (0.64,1.92);
        \draw[very thick, Goldenrod!70!black] (-0.05,0.27) rectangle (0.69,0.69);
        \draw[thick, dashed] (-0.05,0.91) rectangle (0.69,1.33);
        \node[font=\scriptsize] at (0.32,2.4) {$f(\mA_2,\mX_2)$};
      \end{scope}
    \end{scope}
  \end{tikzpicture}
  \caption{One graph, two node orderings (slides 21--28, 45--46). Colours =
  physical nodes; blue cells of $\mA_i$ are 1s. Right: node embeddings
  ($=\mH_1,\mH_2$ on slide 45). Brown node: row $A$ vs.\ row $E$ (gold box);
  green node: row $C$ vs.\ row $D$; blue target node of slide 45: row $D$ vs.\
  row $C$ (dashed).}
  \label{fig:lec03-orders}
\end{figure}

\subsection{Permutation equivariance}\label{sec:lec03-perm-eqv}
\slides{26--30}

\begin{itemize}
  \item \textbf{For node representation:} learn $f$ that maps the nodes of $G$
    to a matrix in $\R^{|V|\times d}$ (each node in $V$ is mapped to a $d$-dim
    embedding).
  \item \Cref{fig:lec03-orders}, right: the representation vector of the brown
    node is row $A$ of $f(\mA_1,\mX_1)$ and row $E$ of $f(\mA_2,\mX_2)$; the green
    node is row $C$ resp.\ row $D$. \textbf{For two orderings, the vector of
    the node at the same position in the graph is the same!}
  \item If the output vector of a node at the same position in the graph
    remains unchanged for any ordering, $f$ is \textbf{permutation
    equivariant}.
\end{itemize}

\begin{definition}[Permutation equivariance]\label{def:lec03-perm-eqv}
  A node function $f:\R^{|V|\times|V|}\times\R^{|V|\times m}\to\R^{|V|\times d}$
  ($f$ maps each node in $V$ to a $d$-dim embedding) is \term{permutation
  equivariant} if
  \begin{equation}\label{eq:lec03-perm-eqv}
    \mP f(\mA,\mX) = f(\mP\mA\mP\tp,\mP\mX)\qquad\text{for any permutation } \mP .
  \end{equation}
\end{definition}

\begin{table}[htbp]
  \centering
  \small
  \begin{tabularx}{\linewidth}{@{}ll>{\raggedright\arraybackslash}X@{}}
    \toprule
    & Condition & Meaning \\
    \midrule
    Perm.-invariant & $f(\mA,\mX)=f(\mP\mA\mP\tp,\mP\mX)$ &
      permute the input, the output stays the same (graph $\to$ vector) \\
    Perm.-equivariant & $\mP f(\mA,\mX)=f(\mP\mA\mP\tp,\mP\mX)$ &
      permute the input, output also permutes accordingly (graph $\to$ matrix) \\
    \bottomrule
  \end{tabularx}
  \caption{Summary: invariance and equivariance (slide 30).}
  \label{tab:lec03-inv-eqv}
\end{table}

\textbf{Examples} (slide 30):
\begin{align}
  f(\mA,\mX)&=\vone\tp\mX: &&\text{invariant:}\quad f(\mP\mA\mP\tp,\mP\mX)=\vone\tp\mP\mX=\vone\tp\mX=f(\mA,\mX)\notag\\
  f(\mA,\mX)&=\mX: &&\text{equivariant:}\quad f(\mP\mA\mP\tp,\mP\mX)=\mP\mX=\mP f(\mA,\mX)\label{eq:lec03-inv-eqv-ex}\\
  f(\mA,\mX)&=\mA\mX: &&\text{equivariant:}\quad f(\mP\mA\mP\tp,\mP\mX)=\mP\mA\mP\tp\mP\mX=\mP\mA\mX=\mP f(\mA,\mX)\notag
\end{align}

\subsection{GNN overview}\label{sec:lec03-overview}
\slides{31--34}

\begin{keyidea}[Bronstein, ICLR 2021 keynote]
  Graph neural networks consist of multiple permutation equivariant /
  invariant functions: permutation-equivariant layer $\to$
  permutation-equivariant layer $\to\cdots\to$ permutation-invariant layer
  (one output vector for the whole graph).
\end{keyidea}
\begin{itemize}
  \item Are other neural network architectures, e.g.\ MLPs, permutation
    invariant / equivariant? \textbf{No.} Switching the order of the input
    leads to different outputs: one MLP gives $-0.6$ for input
    $(-0.53,\,0.15,\,-0.22)$, $-0.8$ for $(0.15,\,-0.53,\,-0.22)$ and $0.02$ for
    $(-0.53,\,-0.22,\,0.15)$.
  \item This explains why \textbf{the naive MLP approach fails for graphs}
    (\cref{fig:lec03-naive}).
\end{itemize}
\begin{keyidea}[Next]
  Design graph neural networks that are permutation invariant / equivariant
  by \textbf{passing and aggregating information from neighbours}.
\end{keyidea}

% =====================================================================
\section{Graph convolutional networks}\label{sec:lec03-gcn}

\subsection{Computation graphs and neighbourhood aggregation}\label{sec:lec03-compgraph}
\slides{35--42}

\begin{itemize}
  \item \textbf{GCN} [Kipf \& Welling, ICLR 2017]. \textbf{Idea:} a node's
    neighbourhood defines a \term{computation graph}: (1) determine the node
    computation graph ($k=1$ and $k=2$ hops around node $i$); (2) propagate and
    transform information (aggregator$_2$: 2-hop $\to$ 1-hop neighbours,
    aggregator$_1$: 1-hop neighbours $\to i$). \textbf{Learn how to propagate
    information across the graph to compute node features.}
  \item \textbf{Key idea:} generate node embeddings based on \textbf{local
    network neighbourhoods} (\cref{fig:lec03-compgraph}, target node $A$).
  \item \textbf{Intuition:} nodes aggregate information from their neighbours
    using \textbf{neural networks} (every box is a neural network); network
    neighbourhood defines a computation graph: \textbf{every node defines a
    computation graph based on its neighbourhood} (\cref{fig:lec03-six-trees}).
  \item \textbf{Deep model: many layers.} Model can be of arbitrary depth:
    nodes have embeddings at each layer; the layer-0 embedding of node $v$ is
    its input feature $\vx_v$; the layer-$k$ embedding gets information from
    nodes that are $k$ hops away.
  \item \textbf{Neighbourhood aggregation:} key distinctions are in how
    different approaches aggregate information across the layers (\emph{what
    is in the box?}).
  \item \textbf{Basic approach:} (1) average messages from neighbours, (2)
    apply a neural network.
\end{itemize}

\begin{figure}[htbp]
  \centering
  \begin{tikzpicture}[scale=0.8, transform shape, lv/.style={font=\small\sffamily}]
    \begin{scope}[shift={(-5.6,-0.9)}, scale=0.85]
      \lecIIIsixgraph{lec03m}{ig}{A}{B}{C}{D}{E}{F}
      \node[font=\scriptsize\sffamily] at (1.8,-0.6) {INPUT GRAPH};
      \draw[arrow] (0.9,2.45) node[above, font=\scriptsize\sffamily] {target node} -- (0.9,1.65);
    \end{scope}
    \node[lec03m, lec03cA] (zA) at (0,0) {A};
    \node[draw, fill=black!10, minimum size=10mm] (box) at (1.4,0) {};
    \node[lec03m, lec03cB] (B) at (3.4,1.72) {B};
    \node[lec03m, lec03cC] (C) at (3.4,0) {C};
    \node[lec03m, lec03cD] (D) at (3.4,-1.65) {D};
    \node[lec03bxd, minimum size=4mm] (bB) at (4.4,1.72) {};
    \node[lec03bxd, minimum size=4mm] (bC) at (4.4,0) {};
    \node[lec03bxd, minimum size=4mm] (bD) at (4.4,-1.65) {};
    \foreach \n/\c/\y in {bA/A/2.0, bC2/C/1.44, cA/A/0.84, cB/B/0.28, cE/E/-0.28, cF/F/-0.84, dA/A/-1.65}
      \node[lec03m, lec03c\c] (\n) at (5.9,\y) {\c};
    \draw[arrow] (box) -- (zA);
    \draw[arrow, dotted] (B) -- (box);
    \draw[arrow, dotted] (C) -- (box);
    \draw[arrow, dotted] (D) -- (box);
    \draw[arrow] (bB) -- (B); \draw[arrow] (bC) -- (C); \draw[arrow] (bD) -- (D);
    \foreach \n in {bA,bC2} \draw[arrow, dotted, black!60] (\n) -- (bB);
    \foreach \n in {cA,cB,cE,cF} \draw[arrow, dotted, black!60] (\n) -- (bC);
    \draw[arrow, dotted, black!60] (dA) -- (bD);
    \foreach \n/\c in {bA/A, bC2/C, cA/A, cB/B, cE/E, cF/F, dA/A}
      \node[right=1pt of \n, font=\small] {$\vx_{\c}$};
    \node[lv] at (1.4,0.9) {Layer-2};
    \node[lv] at (3.4,2.3) {Layer-1};
    \node[lv] at (6.15,2.55) {Layer-0};
  \end{tikzpicture}
  \caption{Computation graph of target node $A$ (slides 37--42): boxes = neural
  networks (``what is in the box?''); layer-0 leaves carry $\vx_u$.}
  \label{fig:lec03-compgraph}
\end{figure}

\begin{figure}[htbp]
  \centering
  \begin{tikzpicture}[scale=0.8, transform shape]
    \node[lec03s, lec03cA] (tsixAr) at (1.584,0.000) {A};
    \node[lec03bxl] (tsixArb) at (1.584,-0.475) {};
    \draw[lec03msg] (tsixArb) -- (tsixAr);
    \node[lec03s, lec03cB] (tsixAc0) at (0.215,-0.950) {B};
    \node[lec03bxd] (tsixAc0b) at (0.215,-1.425) {};
    \draw[lec03msg] (tsixAc0) -- (tsixArb);
    \draw[lec03msg] (tsixAc0b) -- (tsixAc0);
    \node[lec03s, lec03cA] (tsixAc0l0) at (0.000,-1.900) {A};
    \draw[lec03msg] (tsixAc0l0) -- (tsixAc0b);
    \node[lec03s, lec03cC] (tsixAc0l1) at (0.430,-1.900) {C};
    \draw[lec03msg] (tsixAc0l1) -- (tsixAc0b);
    \node[lec03s, lec03cC] (tsixAc1) at (1.655,-0.950) {C};
    \node[lec03bxd] (tsixAc1b) at (1.655,-1.425) {};
    \draw[lec03msg] (tsixAc1) -- (tsixArb);
    \draw[lec03msg] (tsixAc1b) -- (tsixAc1);
    \node[lec03s, lec03cA] (tsixAc1l0) at (1.010,-1.900) {A};
    \draw[lec03msg] (tsixAc1l0) -- (tsixAc1b);
    \node[lec03s, lec03cB] (tsixAc1l1) at (1.440,-1.900) {B};
    \draw[lec03msg] (tsixAc1l1) -- (tsixAc1b);
    \node[lec03s, lec03cE] (tsixAc1l2) at (1.870,-1.900) {E};
    \draw[lec03msg] (tsixAc1l2) -- (tsixAc1b);
    \node[lec03s, lec03cF] (tsixAc1l3) at (2.300,-1.900) {F};
    \draw[lec03msg] (tsixAc1l3) -- (tsixAc1b);
    \node[lec03s, lec03cD] (tsixAc2) at (2.881,-0.950) {D};
    \node[lec03bxd] (tsixAc2b) at (2.881,-1.425) {};
    \draw[lec03msg] (tsixAc2) -- (tsixArb);
    \draw[lec03msg] (tsixAc2b) -- (tsixAc2);
    \node[lec03s, lec03cA] (tsixAc2l0) at (2.881,-1.900) {A};
    \draw[lec03msg] (tsixAc2l0) -- (tsixAc2b);
    \node[lec03s, lec03cB] (tsixBr) at (5.039,0.000) {B};
    \node[lec03bxl] (tsixBrb) at (5.039,-0.475) {};
    \draw[lec03msg] (tsixBrb) -- (tsixBr);
    \node[lec03s, lec03cA] (tsixBc0) at (4.211,-0.950) {A};
    \node[lec03bxd] (tsixBc0b) at (4.211,-1.425) {};
    \draw[lec03msg] (tsixBc0) -- (tsixBrb);
    \draw[lec03msg] (tsixBc0b) -- (tsixBc0);
    \node[lec03s, lec03cB] (tsixBc0l0) at (3.781,-1.900) {B};
    \draw[lec03msg] (tsixBc0l0) -- (tsixBc0b);
    \node[lec03s, lec03cC] (tsixBc0l1) at (4.211,-1.900) {C};
    \draw[lec03msg] (tsixBc0l1) -- (tsixBc0b);
    \node[lec03s, lec03cD] (tsixBc0l2) at (4.641,-1.900) {D};
    \draw[lec03msg] (tsixBc0l2) -- (tsixBc0b);
    \node[lec03s, lec03cC] (tsixBc1) at (5.866,-0.950) {C};
    \node[lec03bxd] (tsixBc1b) at (5.866,-1.425) {};
    \draw[lec03msg] (tsixBc1) -- (tsixBrb);
    \draw[lec03msg] (tsixBc1b) -- (tsixBc1);
    \node[lec03s, lec03cA] (tsixBc1l0) at (5.221,-1.900) {A};
    \draw[lec03msg] (tsixBc1l0) -- (tsixBc1b);
    \node[lec03s, lec03cB] (tsixBc1l1) at (5.651,-1.900) {B};
    \draw[lec03msg] (tsixBc1l1) -- (tsixBc1b);
    \node[lec03s, lec03cE] (tsixBc1l2) at (6.081,-1.900) {E};
    \draw[lec03msg] (tsixBc1l2) -- (tsixBc1b);
    \node[lec03s, lec03cF] (tsixBc1l3) at (6.511,-1.900) {F};
    \draw[lec03msg] (tsixBc1l3) -- (tsixBc1b);
    \node[lec03s, lec03cC] (tsixCr) at (9.518,0.000) {C};
    \node[lec03bxl] (tsixCrb) at (9.518,-0.475) {};
    \draw[lec03msg] (tsixCrb) -- (tsixCr);
    \node[lec03s, lec03cA] (tsixCc0) at (7.841,-0.950) {A};
    \node[lec03bxd] (tsixCc0b) at (7.841,-1.425) {};
    \draw[lec03msg] (tsixCc0) -- (tsixCrb);
    \draw[lec03msg] (tsixCc0b) -- (tsixCc0);
    \node[lec03s, lec03cB] (tsixCc0l0) at (7.411,-1.900) {B};
    \draw[lec03msg] (tsixCc0l0) -- (tsixCc0b);
    \node[lec03s, lec03cC] (tsixCc0l1) at (7.841,-1.900) {C};
    \draw[lec03msg] (tsixCc0l1) -- (tsixCc0b);
    \node[lec03s, lec03cD] (tsixCc0l2) at (8.271,-1.900) {D};
    \draw[lec03msg] (tsixCc0l2) -- (tsixCc0b);
    \node[lec03s, lec03cB] (tsixCc1) at (9.067,-0.950) {B};
    \node[lec03bxd] (tsixCc1b) at (9.067,-1.425) {};
    \draw[lec03msg] (tsixCc1) -- (tsixCrb);
    \draw[lec03msg] (tsixCc1b) -- (tsixCc1);
    \node[lec03s, lec03cA] (tsixCc1l0) at (8.852,-1.900) {A};
    \draw[lec03msg] (tsixCc1l0) -- (tsixCc1b);
    \node[lec03s, lec03cC] (tsixCc1l1) at (9.282,-1.900) {C};
    \draw[lec03msg] (tsixCc1l1) -- (tsixCc1b);
    \node[lec03s, lec03cE] (tsixCc2) at (10.077,-0.950) {E};
    \node[lec03bxd] (tsixCc2b) at (10.077,-1.425) {};
    \draw[lec03msg] (tsixCc2) -- (tsixCrb);
    \draw[lec03msg] (tsixCc2b) -- (tsixCc2);
    \node[lec03s, lec03cC] (tsixCc2l0) at (9.862,-1.900) {C};
    \draw[lec03msg] (tsixCc2l0) -- (tsixCc2b);
    \node[lec03s, lec03cF] (tsixCc2l1) at (10.292,-1.900) {F};
    \draw[lec03msg] (tsixCc2l1) -- (tsixCc2b);
    \node[lec03s, lec03cF] (tsixCc3) at (11.088,-0.950) {F};
    \node[lec03bxd] (tsixCc3b) at (11.088,-1.425) {};
    \draw[lec03msg] (tsixCc3) -- (tsixCrb);
    \draw[lec03msg] (tsixCc3b) -- (tsixCc3);
    \node[lec03s, lec03cC] (tsixCc3l0) at (10.873,-1.900) {C};
    \draw[lec03msg] (tsixCc3l0) -- (tsixCc3b);
    \node[lec03s, lec03cE] (tsixCc3l1) at (11.303,-1.900) {E};
    \draw[lec03msg] (tsixCc3l1) -- (tsixCc3b);
    \begin{scope}[shift={(2.02,0)}]
    \node[lec03s, lec03cD] (tsixDr) at (0.430,-2.900) {D};
    \node[lec03bxl] (tsixDrb) at (0.430,-3.375) {};
    \draw[lec03msg] (tsixDrb) -- (tsixDr);
    \node[lec03s, lec03cA] (tsixDc0) at (0.430,-3.850) {A};
    \node[lec03bxd] (tsixDc0b) at (0.430,-4.325) {};
    \draw[lec03msg] (tsixDc0) -- (tsixDrb);
    \draw[lec03msg] (tsixDc0b) -- (tsixDc0);
    \node[lec03s, lec03cB] (tsixDc0l0) at (0.000,-4.800) {B};
    \draw[lec03msg] (tsixDc0l0) -- (tsixDc0b);
    \node[lec03s, lec03cC] (tsixDc0l1) at (0.430,-4.800) {C};
    \draw[lec03msg] (tsixDc0l1) -- (tsixDc0b);
    \node[lec03s, lec03cD] (tsixDc0l2) at (0.860,-4.800) {D};
    \draw[lec03msg] (tsixDc0l2) -- (tsixDc0b);
    \node[lec03s, lec03cE] (tsixEr) at (3.125,-2.900) {E};
    \node[lec03bxl] (tsixErb) at (3.125,-3.375) {};
    \draw[lec03msg] (tsixErb) -- (tsixEr);
    \node[lec03s, lec03cC] (tsixEc0) at (2.405,-3.850) {C};
    \node[lec03bxd] (tsixEc0b) at (2.405,-4.325) {};
    \draw[lec03msg] (tsixEc0) -- (tsixErb);
    \draw[lec03msg] (tsixEc0b) -- (tsixEc0);
    \node[lec03s, lec03cA] (tsixEc0l0) at (1.760,-4.800) {A};
    \draw[lec03msg] (tsixEc0l0) -- (tsixEc0b);
    \node[lec03s, lec03cB] (tsixEc0l1) at (2.190,-4.800) {B};
    \draw[lec03msg] (tsixEc0l1) -- (tsixEc0b);
    \node[lec03s, lec03cE] (tsixEc0l2) at (2.620,-4.800) {E};
    \draw[lec03msg] (tsixEc0l2) -- (tsixEc0b);
    \node[lec03s, lec03cF] (tsixEc0l3) at (3.050,-4.800) {F};
    \draw[lec03msg] (tsixEc0l3) -- (tsixEc0b);
    \node[lec03s, lec03cF] (tsixEc1) at (3.846,-3.850) {F};
    \node[lec03bxd] (tsixEc1b) at (3.846,-4.325) {};
    \draw[lec03msg] (tsixEc1) -- (tsixErb);
    \draw[lec03msg] (tsixEc1b) -- (tsixEc1);
    \node[lec03s, lec03cC] (tsixEc1l0) at (3.631,-4.800) {C};
    \draw[lec03msg] (tsixEc1l0) -- (tsixEc1b);
    \node[lec03s, lec03cE] (tsixEc1l1) at (4.061,-4.800) {E};
    \draw[lec03msg] (tsixEc1l1) -- (tsixEc1b);
    \node[lec03s, lec03cF] (tsixFr) at (6.326,-2.900) {F};
    \node[lec03bxl] (tsixFrb) at (6.326,-3.375) {};
    \draw[lec03msg] (tsixFrb) -- (tsixFr);
    \node[lec03s, lec03cC] (tsixFc0) at (5.606,-3.850) {C};
    \node[lec03bxd] (tsixFc0b) at (5.606,-4.325) {};
    \draw[lec03msg] (tsixFc0) -- (tsixFrb);
    \draw[lec03msg] (tsixFc0b) -- (tsixFc0);
    \node[lec03s, lec03cA] (tsixFc0l0) at (4.961,-4.800) {A};
    \draw[lec03msg] (tsixFc0l0) -- (tsixFc0b);
    \node[lec03s, lec03cB] (tsixFc0l1) at (5.391,-4.800) {B};
    \draw[lec03msg] (tsixFc0l1) -- (tsixFc0b);
    \node[lec03s, lec03cE] (tsixFc0l2) at (5.821,-4.800) {E};
    \draw[lec03msg] (tsixFc0l2) -- (tsixFc0b);
    \node[lec03s, lec03cF] (tsixFc0l3) at (6.251,-4.800) {F};
    \draw[lec03msg] (tsixFc0l3) -- (tsixFc0b);
    \node[lec03s, lec03cE] (tsixFc1) at (7.046,-3.850) {E};
    \node[lec03bxd] (tsixFc1b) at (7.046,-4.325) {};
    \draw[lec03msg] (tsixFc1) -- (tsixFrb);
    \draw[lec03msg] (tsixFc1b) -- (tsixFc1);
    \node[lec03s, lec03cC] (tsixFc1l0) at (6.831,-4.800) {C};
    \draw[lec03msg] (tsixFc1l0) -- (tsixFc1b);
    \node[lec03s, lec03cF] (tsixFc1l1) at (7.261,-4.800) {F};
    \draw[lec03msg] (tsixFc1l1) -- (tsixFc1b);
    \end{scope}
    \draw[rounded corners=4pt, thick, WildStrawberry!80!black] (-0.3,0.35) rectangle (11.6,-2.15);
    \draw[rounded corners=4pt, thick, RoyalBlue!80!black] (1.72,-2.55) rectangle (9.58,-5.05);
    \node[font=\scriptsize, WildStrawberry!80!black, anchor=south west] at (-0.3,0.37) {(3) train on a batch of compute graphs ($A,B,C$)};
    \node[font=\scriptsize, RoyalBlue!80!black, anchor=north west] at (1.72,-5.07) {(4) even for nodes we never trained on ($D,E,F$)};
  \end{tikzpicture}
  \caption{Every node defines its own computation graph (slide 39); reused on
  slides 56--57 for training and for generating embeddings as needed.}
  \label{fig:lec03-six-trees}
\end{figure}

\subsection{The math: deep encoder of a GCN}\label{sec:lec03-gcn-math}
\slides{43, 48}

\textbf{Basic approach:} average neighbour messages and apply a neural network.

\begin{definition}[GCN deep encoder, GraphSAGE-style mean aggregation (Hamilton et al., 2017) ]\label{def:lec03-gcn}
  \begin{align}
    \vh_v^{(0)} &= \vx_v, \label{eq:lec03-gcn-init}\\
    \vh_v^{(k+1)} &= \sigma\Bigl(\mW_k\sum_{u\in\Nb(v)}\frac{\vh_u^{(k)}}{|\Nb(v)|}
      + \mB_k\vh_v^{(k)}\Bigr),\qquad \forall k\in\{0,\dots,K-1\},
      \label{eq:lec03-gcn-update}\\
    \vz_v &= \vh_v^{(K)}. \label{eq:lec03-gcn-out}
  \end{align}
\end{definition}
\begin{itemize}
  \item $\vh_v^{(0)}=\vx_v$: initial 0-th layer embeddings = node features;
    $\vh_v^{(k)}$: hidden representation of $v$ at layer $k$;
    $\sum_{u\in\Nb(v)}\vh_u^{(k)}/|\Nb(v)|$: \textbf{average of neighbours'
    previous-layer embeddings} (summation is a permutation invariant function);
    $\sigma$: non-linearity (e.g.\ ReLU); $K$: total number of layers; $\vz_v$:
    embedding after $K$ layers of neighbourhood aggregation (final embedding).
  \item \textbf{Model parameters} (slide 48): trainable weight matrices (what we
    learn): $\mW_k$ = weight matrix for \textbf{neighbourhood aggregation},
    $\mB_k$ = weight matrix for \textbf{transforming the hidden vector of
    self}. We can feed these \textbf{embeddings into any loss function} and run
    SGD to \textbf{train the weight parameters}.
\end{itemize}

\subsection{GCN: invariance and equivariance}\label{sec:lec03-gcn-inv}
\slides{44--46}

\begin{itemize}
  \item \textbf{Given a node, the GCN that computes its embedding is
    permutation invariant.} Computation graph of target node $D$:
    $D \leftarrow [\,\cdot\,] \leftarrow A \leftarrow [\,\cdot\,] \leftarrow
    \{B,C,D\}$ with \textbf{shared} NN weights; the \textbf{average} of
    neighbours' previous-layer embeddings is permutation invariant.
  \item \textbf{Considering all nodes in a graph, GCN computation is
    permutation equivariant}: permute the input, the output also permutes
    accordingly (\cref{fig:lec03-orders}: target node $D$ in order~1 = $C$ in
    order~2).
  \item \textbf{Detailed reasoning:} (1) the rows of \textbf{input node
    features} and \textbf{output embeddings} are \textbf{aligned}; (2)
    computing the embedding of \textbf{a given node} with GCN is
    \textbf{invariant}; (3) so, after permutation, the \textbf{location} of a
    given node in the input feature matrix is changed, and \textbf{the output
    embedding of a given node stays the same} (the colours of node feature and
    embedding are matched) $\Rightarrow$ permutation equivariant.
\end{itemize}

\subsection{Matrix formulation}\label{sec:lec03-matrix}
\slides{49--50}

\textbf{Many aggregations can be performed efficiently by (sparse) matrix
operations.} Let $\mH^{(k)}$ be the matrix of hidden embeddings (row $i$ =
$\vh_i^{(k)}$) and $\mD$ the diagonal degree matrix:
\begin{gather}
  \mH^{(k)} = [\vh_1^{(k)}\cdots\vh_{|V|}^{(k)}]\tp,\qquad
  \sum_{u\in\Nb_v}\vh_u^{(k)} = \mA_{v,:}\mH^{(k)}, \label{eq:lec03-H}\\
  D_{v,v}=\mathrm{Deg}(v)=|\Nb(v)|,\qquad
  D^{-1}_{v,v}=1/|\Nb(v)|\quad(\mD^{-1}\text{ is also diagonal}). \label{eq:lec03-D}
\end{gather}
Therefore
\begin{equation}\label{eq:lec03-mean-matrix}
  \sum_{u\in\Nb(v)}\frac{\vh_u^{(k)}}{|\Nb(v)|}\ \Longrightarrow\ \mD^{-1}\mA\mH^{(k)} .
\end{equation}
\begin{caution}
  Slide 49 writes $\sum_u\vh_u^{(k-1)}/|\Nb(v)| \Rightarrow \mH^{(k+1)}=\mD^{-1}\mA\mH^{(k)}$:
  indices mixed; $\mD^{-1}\mA\mH^{(k)}$ is only the neighbour mean, the full layer is \cref{eq:lec03-gcn-matrix}.
\end{caution}
Re-writing the update function \cref{eq:lec03-gcn-update} in matrix form:
\begin{equation}\label{eq:lec03-gcn-matrix}
  \mH^{(k+1)} = \sigma\bigl(\textcolor{cautionred}{\tilde\mA\mH^{(k)}\mW_k\tp}
  + \textcolor{keyblue}{\mH^{(k)}\mB_k\tp}\bigr),
  \qquad \tilde\mA=\mD^{-1}\mA
\end{equation}
\begin{itemize}
  \item \textcolor{cautionred}{Red: neighbourhood aggregation};
    \textcolor{keyblue}{blue: self transformation} (picture: red edges from
    the neighbours into $v$, blue self-loop). In practice, efficient
    \textbf{sparse matrix multiplication} can be used ($\tilde\mA$ is sparse).
  \item \textbf{Note:} not all GNNs can be expressed in a simple matrix form,
    when the aggregation function is complex.
\end{itemize}

\subsection{Training the model}\label{sec:lec03-training}
\slides{47, 51--54}

\begin{itemize}
  \item How do we train the GCN to generate embeddings (e.g.\ $\vz_A$)?
    \textbf{Need to define a loss function on the embeddings.} Node embedding
    $\vz_v$ is a function of the input graph.
  \item \textbf{Supervised setting:} minimise the loss $\Loss$ ($\vy$: node
    label; $\Loss$ could be L2 if $\vy$ is a real number, or cross entropy if
    $\vy$ is categorical):
\end{itemize}
\begin{equation}\label{eq:lec03-supervised}
  \min_\Theta\ \Loss\bigl(\vy, f_\Theta(\vz_v)\bigr)
\end{equation}
\begin{itemize}
  \item \textbf{Unsupervised setting:} no node label available; \textbf{use the
    graph structure as the supervision!} One possible idea: ``similar'' nodes
    have similar embeddings:
\end{itemize}
\begin{gather}
  \min_\Theta\ \Loss = \sum_{z_u,z_v}\mathrm{CE}\bigl(y_{u,v},\DEC(\vz_u,\vz_v)\bigr),
  \label{eq:lec03-unsupervised}\\
  \mathrm{CE}\bigl(\vy,f(\vx)\bigr) = -\sum_{i=1}^{C} y_i\log f_\Theta(\vx)_i, \qquad C = \textbf{num classes, binary here} \label{eq:lec03-ce-unsup}
\end{gather}
\begin{itemize}
  \item $y_{u,v}=1$ when nodes $u$ and $v$ are \textbf{similar};
    $\vz_u=f_\Theta(u)$; $\DEC(\cdot,\cdot)$ is the dot product. CE: $y_i$ and
    $f_\Theta(\vx)_i$ are the \textbf{actual} and \textbf{predicted} values of
    the $i$-th class; intuition: the lower the loss, the closer the prediction
    is to one-hot.
  \item \textbf{Node similarity} can be anything from Lecture~2, e.g.\ a loss
    based on \textbf{random walks} (node2vec, DeepWalk, struc2vec) or
    \textbf{matrix factorisation}.
\end{itemize}
\begin{caution}
  Slide 52 takes $\DEC$ = dot product inside a cross entropy; CE needs a
  probability, so read $\DEC(\vz_u,\vz_v)=\sigma(\vz_u\tp\vz_v)$.
\end{caution}
\begin{itemize}
  \item \textbf{Supervised training:} directly train the model for a supervised
    task (e.g.\ \textbf{node classification}): \emph{safe or toxic drug?} for
    node $A$ in a drug--drug interaction network. Use cross entropy loss
    ($\vz_v$: encoder output = node embedding; $\vtheta$: classification
    weights; $y_v$: node class label):
\end{itemize}
\begin{equation}\label{eq:lec03-bce}
  \Loss = -\sum_{v\in V}\Bigl[\,y_v\log\bigl(\sigma(\vz_v\tp\vtheta)\bigr)
  + (1-y_v)\log\bigl(1-\sigma(\vz_v\tp\vtheta)\bigr)\Bigr]
\end{equation}
\begin{caution}
  Slide 54 omits the outer brackets (literally $-\sum_v$ applies to the first
  term only); both terms belong inside $-\sum_v[\dots]$, as in \cref{eq:lec03-bce}.
\end{caution}

\subsection{Model design and inductive capability}\label{sec:lec03-inductive}
\slides{55--60}

\begin{itemize}
  \item \textbf{Model design overview:} (1) define a neighbourhood aggregation
    function; (2) define a loss function on the embeddings; (3) train on a set
    of nodes, i.e.\ a batch of compute graphs; (4) generate embeddings for
    nodes as needed, \textbf{even for nodes we never trained on!}
    (\cref{fig:lec03-six-trees}).
  \item \textbf{Inductive capability: the same aggregation parameters are
    shared for all nodes} (compute graphs of $A$ and $B$ share $\mW_k,\mB_k$ at
    each level) $\Rightarrow$ the number of model parameters is sublinear in
    $|V|$ and we can \textbf{generalise to unseen nodes}!
  \item \textbf{New graphs:} train on one graph, generalise to a new graph
    (inductive node embedding $\to$ generalise to entirely unseen graphs); e.g.\
    train on a protein interaction graph from model organism $A$ and generate
    embeddings on newly collected data about organism $B$.
  \item \textbf{New nodes:} train with snapshot $\to$ new node arrives $\to$
    generate embedding $\vz_u$ for the new node. Many application settings
    constantly encounter previously unseen nodes (e.g.\ Reddit, YouTube,
    Google Scholar) $\Rightarrow$ need to generate new embeddings ``on the
    fly''.
\end{itemize}

\subsection{Interactive visualisation: GNN 101}\label{sec:lec03-gnn101}
\slides{61--62}

\begin{itemize}
  \item \textbf{GNN 101} (\url{https://visual-intelligence-umn.github.io/GNN-101};
    \emph{Visual Learning of GNNs in Your Web Browser}): inner working of
    a GNN with a \textbf{hierarchical level of detail}; level one = layers
    (input $\to$ GCNConv1--3 $\to$ global mean pooling $\to$ output), level two
    = one node in one layer (vector summation, $64\times64$ weight matrix, bias,
    ReLU), level three = the matrix multiplication. Parameters shown as
    heatmaps, curves/animations = computation process; hovering over terms in
    the equation highlights the corresponding visualisation.
  \item Example: GCN graph classification on \textbf{MUTAG} (graph = chemical
    compound, nodes = atoms, edges = bonds; mutagenic on \emph{Salmonella
    typhimurium} or not; 188 graphs, 150 train, 38 test).
  \item \textbf{Other functions:} GNN variants and tasks: GAT (attention score,
    e.g.\ $\alpha_{497,194}=0.06$, softmax of
    $e_{i,j}=\mathrm{LeakyReLU}(\va\mW\vx_i+\va\mW\vx_j)$), node classification,
    edge prediction (e.g.\ $0.88\to$ True); e.g.\ GAT link prediction on
    \textbf{Twitch} (nodes = users, edges = mutual follower relationships;
    77{,}774 edges, 60{,}052 in training). Toggle node-link subgraph view /
    adjacency \textbf{matrix view}. Tool formulas (GCN; GAT):
\end{itemize}
\begin{equation*}
  \vx_i'=\sigma\Bigl(\mW\!\!\sum_{j\in\Nb(i)\cup\{i\}}\!\frac{e_{j,i}}{\sqrt{\hat d_j\hat d_i}}\vx_j+\vb\Bigr),
  \qquad
  \vx_i'=\sigma\Bigl(\mW\!\!\sum_{j\in\Nb(i)\cup\{i\}}\!\alpha_{i,j}\vx_j+\vb\Bigr)
\end{equation*}

% =====================================================================
\section{GNNs subsume CNNs (and transformers)}\label{sec:lec03-subsume}

\subsection{CNN as a special GNN}\label{sec:lec03-cnn}
\slides{63--68}

\begin{itemize}
  \item \textbf{Architecture comparison:} how do GNNs compare to prominent
    architectures such as convolutional neural nets?
  \item CNN layer with $3\times3$ filter (example: image patch
    $\bigl(\begin{smallmatrix}9&4&1\\1&1&1\\1&2&1\end{smallmatrix}\bigr)$, CNN
    weights $\bigl(\begin{smallmatrix}0&2&1\\4&1&0\\1&0&1\end{smallmatrix}\bigr)$
    $\Rightarrow$ output $16$):
\end{itemize}
\begin{equation}\label{eq:lec03-cnn}
  \text{CNN:}\quad \vh_v^{(l+1)} = \sigma\Bigl(\sum_{u\in\Nb(v)\cup\{v\}}\mW_l^{u}\vh_u^{(l)}\Bigr),
  \qquad\forall l\in\{0,\dots,L-1\},
\end{equation}
where \textbf{$\Nb(v)$ represents the 8 neighbour pixels of $v$}. GNN vs.\ CNN
(image vs.\ graph: 8 neighbours plus a self-loop):
\begin{align}
  \text{GNN:}\quad \vh_v^{(l+1)} &= \sigma\Bigl(\textcolor{cautionred}{\mW_l}\sum_{u\in\Nb(v)}\frac{\vh_u^{(l)}}{|\Nb(v)|} + \mB_l\vh_v^{(l)}\Bigr),
  \qquad\forall l\in\{0,\dots,L-1\} \label{eq:lec03-gnn-vs}\\
  \text{CNN (rewritten):}\quad \vh_v^{(l+1)} &= \sigma\Bigl(\sum_{u\in\Nb(v)}\textcolor{cautionred}{\mW_l^{u}}\vh_u^{(l)} + \mB_l\vh_v^{(l)}\Bigr),
  \qquad\forall l\in\{0,\dots,L-1\} \label{eq:lec03-cnn-vs}
\end{align}
\begin{itemize}
  \item \textbf{Key difference:} we can learn different $\mW_l^u$ for different
    ``neighbour'' $u$ for pixel $v$ on the image. The reason is we can pick an
    order for the 9 neighbours using \textbf{relative position} to the centre
    pixel: $\{(-1,-1),(-1,0),(-1,1),\dots,(1,1)\}$.
\end{itemize}
\begin{keyidea}[Slide 68]
  \begin{itemize}
    \item \textbf{CNN can be seen as a special GNN with fixed neighbour size and
      ordering:} the size of the filter is pre-defined for a CNN; the advantage
      of GNN is it processes arbitrary graphs with different degrees for each
      node.
    \item \textbf{CNN is not permutation invariant/equivariant:} switching the
      order of pixels leads to different outputs.
  \end{itemize}
\end{keyidea}

\subsection{Transformer as a special GNN}\label{sec:lec03-transformer}
\slides{69--71}

\begin{itemize}
  \item \textbf{Transformer} [Vaswani et al., NeurIPS 2017]: among the most
    popular architectures; great performance in many sequence modelling tasks
    (diagram: multi-head attention; scaled dot-product attention = MatMul($Q,K$)
    $\to$ scale $\to$ mask (opt.) $\to$ softmax $\to$ MatMul with $V$).
    \textbf{Key component: self-attention:} every token/word attends to all
    the other tokens/words via matrix calculation (``I am a Stanford
    student'').
\end{itemize}
\begin{definition}[Attention]\label{def:lec03-attention}
  Given a set of vector \emph{values} and a vector \emph{query},
  \term{attention} is a technique to compute a weighted sum of the values,
  dependent on the query.
\end{definition}
\begin{itemize}
  \item Each token/word has a \textbf{value vector} (can be seen as the
    representation of the token/word) and a \textbf{query vector} (used to
    calculate the attention scores = weights in the weighted sum). E.g.\
    ``The animal didn't cross the street because it was too tired'': ``it''
    attends to ``animal''; ``\dots too wide'': ``it'' attends to ``street''.
\end{itemize}
\begin{keyidea}[GNN vs.\ transformer]
  Transformer layer can be seen as a \textbf{special GNN that runs on a
  fully-connected ``word'' graph}! Since each word attends to \textbf{all the
  other words}, the computation graph of a transformer layer is identical to
  that of a GNN on the fully-connected ``word'' graph
  (\cref{fig:lec03-transformer}; blog: ``Transformers are Graph Neural
  Networks'', towardsdatascience.com).
\end{keyidea}

\begin{figure}[htbp]
  \centering
  \begin{tikzpicture}[scale=0.8, transform shape]
    \foreach \w/\x [count=\i] in {I/0, am/1.3, a/2.6, Stanford/4.15, student/5.85} {
      \node[ellipse, draw=Goldenrod!70!black, fill=Goldenrod!45, thick, minimum width=11mm,
            minimum height=6.5mm, inner sep=1pt, font=\scriptsize] (t\i) at (\x,0) {\w};
    }
    \foreach \j/\lo in {2/1.2, 3/1.0, 4/0.9, 5/0.85} {
      \draw[-{Stealth[length=1.6mm]}, Goldenrod!70!black, thick] (t1.south) to[out=-70, in=-110, looseness=\lo] (t\j.south);
    }
    \node[font=\small, cautionred] at (2.9,-2.0) {text};
    \begin{scope}[shift={(9.7,0.05)}]
      \foreach \a [count=\i] in {90, 18, -54, -126, 162} {
        \node[circle, draw=keyblue, very thick, minimum size=6mm] (w\i) at (\a:1.35) {};
      }
      \foreach \i in {1,...,5} \foreach \j in {1,...,5} {
        \ifnum\i<\j \draw[thick, black!45] (w\i) -- (w\j); \fi
      }
      \node[above=1pt of w1, font=\small] {I};
      \node[right=1pt of w2, font=\small] {student};
      \node[below right=0pt and 0pt of w3, font=\small] {Stanford};
      \node[below left=0pt and 0pt of w4, font=\small] {a};
      \node[left=1pt of w5, font=\small] {am};
      \node[font=\small, cautionred] at (0,-2.05) {(complete) graph};
    \end{scope}
  \end{tikzpicture}
  \caption{Transformer layer = GNN on the complete ``word'' graph (slide 71);
  only the arrows from ``I'' are drawn.}
  \label{fig:lec03-transformer}
\end{figure}

% =====================================================================
\section{Background: basics of deep learning}\label{sec:lec03-dl}

\subsection{Machine learning as optimisation}\label{sec:lec03-ml-opt}
\slides{73--76}

\begin{itemize}
  \item \textbf{Supervised learning:} we are given input $\vx$, and the goal is
    to predict label $\vy$. Input $\vx$ can be: vectors of real numbers;
    sequences (natural language); matrices (images); graphs (potentially with
    node and edge features).
  \item We formulate the task as an \textbf{optimisation problem}:
\end{itemize}
\begin{equation}\label{eq:lec03-objective}
  \min_{\Theta}\ \Loss\bigl(\vy, f(\vx)\bigr)\qquad\text{(objective function)}
\end{equation}
\begin{itemize}
  \item $\Theta$: a set of \textbf{parameters} we optimise; could contain one or
    more scalars, vectors, matrices, \dots; e.g.\ $\Theta=\{\mZ\}$ in the shallow
    encoder (the embedding lookup). $\Loss$: \textbf{loss function}. Example: L2 loss
    $\Loss(\vy,f(\vx))=\norm{\vy-f(\vx)}_2$. Other common loss functions: L1
    loss, Huber loss, max margin (hinge loss), cross entropy, \dots{} (see
    \url{https://pytorch.org/docs/stable/nn.html\#loss-functions}).
  \item \textbf{Loss function example:} one common loss for classification is
    \textbf{cross entropy (CE)}. Label $\vy$ is a categorical vector
    (\textbf{one-hot encoding}), e.g.\ $\vy=[0,0,1,0,0]$ ($\vy$ is of class
    ``3''). $f(\vx)=\softmax(g(\vx))$, e.g.\ $f(\vx)=[0.1,0.3,0.4,0.1,0.1]$, with
    ($C$ = number of classes, $g(\vx)_i$ = $i$-th coordinate of the vector
    output of $g$):
\end{itemize}
\begin{equation}\label{eq:lec03-softmax}
  f(\vx)_i = \frac{e^{g(\vx)_i}}{\sum_{j=1}^{C} e^{g(\vx)_j}}
\end{equation}
\begin{caution}
  Slide 76 says ``Recall from lecture 3'' for the softmax: stale
  cross-reference (this is Lecture~3).
\end{caution}
\begin{equation}\label{eq:lec03-ce}
  \mathrm{CE}\bigl(\vy,f(\vx)\bigr) = -\sum_{i=1}^{C} y_i\log f(\vx)_i ,
  \qquad
  \Loss = \sum_{(\vx,\vy)\in\calT}\mathrm{CE}\bigl(\vy,f(\vx)\bigr)
\end{equation}
\begin{itemize}
  \item $y_i$, $f(\vx)_i$: \textbf{actual} and \textbf{predicted} values of the
    $i$-th class; intuition: the lower the loss, the closer the prediction is
    to one-hot. \textbf{Total loss over all training examples:} $\calT$ =
    training set containing all pairs of data and labels $(\vx,\vy)$.
\end{itemize}

\subsection{Gradient descent and SGD}\label{sec:lec03-sgd}
\slides{77--80}

\begin{itemize}
  \item How to optimise the objective function? \textbf{Gradient vector:}
    direction and rate of fastest increase ($\Theta_1,\Theta_2,\dots$:
    components of $\Theta$; entries = partial derivatives):
\end{itemize}
\begin{equation}\label{eq:lec03-grad}
  \nabla_\Theta\Loss = \Bigl(\frac{\partial\Loss}{\partial\Theta_1},\frac{\partial\Loss}{\partial\Theta_2},\dots\Bigr)
\end{equation}
\begin{itemize}
  \item Recall \textbf{directional derivative} of a multi-variable function
    (e.g.\ $\Loss$) along a given vector = instantaneous rate of change of the
    function along the vector. \textbf{Gradient is the directional derivative
    in the direction of largest increase.}
  \item \textbf{Gradient descent} (iterative algorithm): repeatedly update
    weights in the (opposite) direction of gradients until convergence:
\end{itemize}
\begin{equation}\label{eq:lec03-gd}
  \Theta \leftarrow \Theta - \eta\,\nabla_\Theta\Loss
\end{equation}
\begin{itemize}
  \item \textbf{Training:} optimise $\Theta$ iteratively; \textbf{iteration}: 1
    step of gradient descent. \textbf{Learning rate (LR) $\eta$:}
    hyperparameter that controls the size of the gradient step; can vary over
    the course of training (LR scheduling).
  \item \textbf{Ideal termination condition:} gradient $=\vzero$. In practice,
    stop training if it no longer improves performance on the
    \textbf{validation set} (part of the dataset held out from training).
  \item \textbf{Problem with gradient descent:} the exact gradient requires
    $\nabla_\Theta\Loss(\vy,f(\vx))$ with $\vx$ = the \textbf{entire} dataset
    (sum of gradient contributions over all points; modern datasets often
    contain billions of points) $\Rightarrow$ extremely expensive for every
    step.
  \item \textbf{Solution: stochastic gradient descent (SGD):} at every step,
    pick a different \textbf{minibatch} $\mathcal{B}$ containing a subset of
    the dataset, use it as input $\vx$.
  \item \textbf{Minibatch SGD concepts:} \textbf{batch size} = number of data
    points in a minibatch (e.g.\ number of nodes for a node classification
    task); \textbf{iteration} = 1 step of SGD on a minibatch; \textbf{epoch} =
    one full pass over the dataset (\# iterations = dataset size / batch size).
  \item \textbf{SGD is an unbiased estimator of the full gradient}, but no
    guarantee on the rate of convergence; in practice often requires tuning of
    the learning rate. Common optimisers that improve over SGD: Adam, Adagrad,
    Adadelta, RMSprop, \dots
\end{itemize}

\subsection{Neural network functions and back-propagation}\label{sec:lec03-backprop}
\slides{81--85}

\begin{itemize}
  \item Objective $\min_\Theta\Loss(\vy,f(\vx))$; in deep learning, $f$ can be
    very complex. Start simple: linear function $f(\vx)=\mW\cdot\vx$,
    $\Theta=\{\mW\}$.
    \begin{itemize}
      \item If $f$ returns a scalar, $\mW$ is a learnable \textbf{vector}:
        $\nabla_\mW f=\bigl(\frac{\partial f}{\partial w_1},\frac{\partial f}{\partial w_2},\frac{\partial f}{\partial w_3},\dots\bigr)$.
      \item If $f$ returns a vector, $\mW$ is the \textbf{weight matrix} and
        $\nabla_\mW f$ is the \textbf{Jacobian matrix of $f$}, drawn as
        $\bigl[\begin{smallmatrix}\partial f_1/\partial w_{11} & \partial f_2/\partial w_{12}\\
        \partial f_1/\partial w_{21} & \partial f_2/\partial w_{22}\end{smallmatrix}\bigr]$.
    \end{itemize}
  \item \textbf{Back-propagation intuition:} to minimise $\Loss$ we need to
    evaluate the gradient $\nabla_\mW\Loss$, i.e.\ derive the derivative of
    $\Loss$.
  \item \textbf{Overview of back-propagation:} $\Loss$ is composed from some set
    of predefined building block functions $g(\cdot)$; for each such $g$ we also
    have its derivative $g'$; then we can automatically compute
    $\nabla_\Theta\Loss$ by evaluating the appropriate $g'$ on the minibatch
    $\mathcal{B}$.
\end{itemize}
\begin{caution}
  Slide 81's $2\times2$ ``Jacobian'' is schematic (true entries: $\partial f_i/\partial w_{jk}$).
  Slide 82 writes $\nabla_\mW\Loss=(\frac{\partial f}{\partial w_1},\dots)$; correct: $(\frac{\partial\Loss}{\partial w_1},\frac{\partial\Loss}{\partial w_2},\dots)$.
\end{caution}
\begin{itemize}
  \item More complex function: $f(\vx)=\mW_2(\mW_1\vx)$,
    $\Theta=\{\mW_1,\mW_2\}$; in other words $h(\vx)=\mW_1\vx$,
    $g(\vz)=\mW_2\vz$. Recall the \textbf{chain rule}:
\end{itemize}
\begin{equation}\label{eq:lec03-chain}
  \frac{\mathrm{d}f}{\mathrm{d}x} = \frac{\mathrm{d}g}{\mathrm{d}h}\cdot\frac{\mathrm{d}h}{\mathrm{d}x}
  \quad\text{or}\quad f'(x) = g'\bigl(h(x)\bigr)\,h'(x);
  \qquad\text{example: }
  \nabla_\vx f = \frac{\partial f}{\partial(\mW_1\vx)}\cdot\frac{\partial(\mW_1\vx)}{\partial\vx}
\end{equation}
\begin{itemize}
  \item \textbf{Back-propagation:} use of the \textbf{chain rule} to propagate
    gradients of intermediate steps, and finally obtain the gradient of $\Loss$
    w.r.t.\ $\Theta$.
  \item \textbf{Example:} simple 2-layer linear network ($x_1,x_2\to$ 2 hidden
    units $\to f(\vx)$), $f(\vx)=g(h(\vx))=\mW_2(\mW_1\vx)$. \textbf{Hidden
    layer:} intermediate representation of input $\vx$; here $h(\vx)=\mW_1\vx$
    denotes the hidden layer, $f(\vx)=\mW_2h(\vx)$. The loss sums the L2 loss
    in a minibatch $\mathcal{B}$:
\end{itemize}
\begin{equation}\label{eq:lec03-bp-loss}
  \Loss = \sum_{(\vx,\vy)\in\mathcal{B}}\norm{\vy-f(\vx)}_2
\end{equation}
\begin{itemize}
  \item \textbf{Forward propagation:} compute loss starting from input:
    $\vx \xrightarrow{\times\mW_1} h \xrightarrow{\times\mW_2} g \xrightarrow{\text{loss}}\Loss$.
  \item \textbf{Back-propagation} to compute the gradient of
    $\Theta=\{\mW_1,\mW_2\}$: start from the loss, compute the gradient
    (compute backwards):
\end{itemize}
\begin{equation}\label{eq:lec03-backprop}
  \frac{\partial\Loss}{\partial\mW_2} = \frac{\partial\Loss}{\partial f}\cdot\frac{\partial f}{\partial\mW_2},
  \qquad
  \frac{\partial\Loss}{\partial\mW_1} = \frac{\partial\Loss}{\partial f}\cdot\frac{\partial f}{\partial h}\cdot\frac{\partial h}{\partial\mW_1}
\end{equation}
\begin{caution}
  Slide 84 prints $\norm{(\vy,-f(\vx))}_2$ (comma = typo). Slide 85 writes
  $\frac{\partial\Loss}{\partial\mW_1}=\frac{\partial\Loss}{\partial f}\cdot\frac{\partial f}{\partial\mW_2}\cdot\frac{\partial\mW_2}{\partial\mW_1}$;
  correct: the path goes through $h$ (\cref{eq:lec03-backprop}).
\end{caution}

\subsection{Non-linearity and the multi-layer perceptron}\label{sec:lec03-mlp}
\slides{86--89}

\begin{itemize}
  \item In $f(\vx)=\mW_2(\mW_1\vx)$, $\mW_2\mW_1$ is another matrix (vector, if
    we do binary classification) $\Rightarrow$ $f(\vx)$ is still linear w.r.t.\
    $\vx$ no matter how many weight matrices we compose.
  \item We introduce \textbf{non-linearity}: rectified linear unit (ReLU) and
    sigmoid:
\end{itemize}
\begin{equation}\label{eq:lec03-nonlin}
  \ReLU(x)=\max(x,0),\qquad \sigma(x)=\frac{1}{1+e^{-x}}
\end{equation}
\begin{itemize}
  \item \textbf{Multi-layer perceptron (MLP):} each layer combines linear
    transformation and non-linearity:
\end{itemize}
\begin{equation}\label{eq:lec03-mlp}
  \vx^{(l+1)} = \sigma\bigl(\mW_l\vx^{(l)} + \vb^{l}\bigr)
\end{equation}
\begin{itemize}
  \item $\mW_l$: weight matrix that transforms the hidden representation at
    layer $l$ to layer $l+1$; $\vb^l$: bias at layer $l$, added to the linear
    transformation of $\vx^{(l)}$; $\sigma$: non-linearity function (e.g.\
    sigmoid).
  \item Example: $\vx$ 2-dimensional ($x_1,x_2$) $\to$ 3-dimensional hidden
    representation $\to$ 2 units $\to$ 1-dimensional output; every layer:
    linear transformation + non-linearity.
  \item \textbf{Summary} (slide 88): the recipe of slide 13
    (\cref{eq:lec03-dl-recipe}): objective $\min_\Theta\Loss(\vy,f(\vx))$; $f$ =
    linear layer, MLP or other NN (e.g.\ a GNN); sample a minibatch; forward
    propagation; back-propagation ($\nabla_\mW\Loss$ via chain rule); SGD over
    many iterations. Slide 89: outline, back to deep learning for graphs
    (\cref{sec:lec03-dlg}).
\end{itemize}

% ---------------------------------------------------------------------
\bigskip\phantomsection
\addcontentsline{toc}{section}{Lecture 3 at a glance}
\begin{summarybox}[Lecture 3 at a glance]
  \begin{itemize}
    \item Deck summary (slide 72): deep learning for graphs = multiple layers of
      embedding transformation; each layer takes the previous-layer embedding
      and aggregates neighbours and self. GCN = mean aggregation, matrix form.
      GNN is a general architecture (CNN is a special GNN).
    \item Shallow encoders: $\bigO(|V|d)$ parameters, transductive, no features.
      Graphs have no node ordering $\Rightarrow$ invariance
      $f(\mA,\mX)=f(\mP\mA\mP\tp,\mP\mX)$ / equivariance
      $\mP f(\mA,\mX)=f(\mP\mA\mP\tp,\mP\mX)$; an MLP on $[\mA,\mX]$ is neither.
      GNN = equivariant layers + invariant readout.
    \item GCN:
      $\vh_v^{(k+1)}=\sigma(\mW_k\sum_{u\in\Nb(v)}\vh_u^{(k)}/|\Nb(v)|+\mB_k\vh_v^{(k)})$,
      $\vh_v^{(0)}=\vx_v$, $\vz_v=\vh_v^{(K)}$;
      $\mH^{(k+1)}=\sigma(\tilde\mA\mH^{(k)}\mW_k\tp+\mH^{(k)}\mB_k\tp)$,
      $\tilde\mA=\mD^{-1}\mA$. Per node invariant, all nodes equivariant.
    \item Training: supervised $\min_\Theta\Loss(\vy,f_\Theta(\vz_v))$ (e.g.\
      BCE with $\sigma(\vz_v\tp\vtheta)$); unsupervised
      $\sum\mathrm{CE}(y_{u,v},\DEC(\vz_u,\vz_v))$. Shared $\mW_k,\mB_k$
      $\Rightarrow$ inductive (unseen nodes, new graphs).
    \item CNN: own $\mW_l^u$ per relative position, fixed size/order, not
      equivariant. Transformer = GNN on the complete word graph.
    \item DL basics: $\Theta\leftarrow\Theta-\eta\nabla_\Theta\Loss$ (minibatch
      SGD); back-prop = chain rule; MLP $\vx^{(l+1)}=\sigma(\mW_l\vx^{(l)}+\vb^l)$.
  \end{itemize}
\end{summarybox}
